\title{Large Language Models Can Automatically Engineer Features for Few-Shot Tabular Learning}

\begin{abstract}
Large Language Models (LLMs) have demonstrated impressive proficiency in solving complex and novel reasoning challenges, making them highly promising tools for tabular learning—a critical need across many practical domains. In this work, we introduce an innovative in-context learning paradigm called \textsf{FeatLLM}, in which LLMs act as feature creators, generating an enhanced input dataset specifically tailored for tabular prediction tasks. These engineered features are leveraged by a basic downstream machine learning model, such as linear regression, which enables robust few-shot learning and high predictive performance. During inference, \textsf{FeatLLM} relies solely on the trained simple predictive model and the features formulated during training. This distinguishes \textsf{FeatLLM} from prior LLM-based methods, as it obviates the necessity of querying the LLM for every inference instance. Additionally, it functions effectively using only API access to LLMs, thereby overcoming common prompt length restrictions. Empirical evaluations on an array of tabular datasets from diverse fields demonstrate that \textsf{FeatLLM} discovers high-quality rules, achieving a significant performance improvement—by an average of 10\%—over methods like TabLLM and STUNT.
\end{abstract}

\section{Introduction}
Large Language Models (LLMs) have, in recent years, exhibited a remarkable capacity to address novel tasks with minimal or no need for task-specific retraining~\cite{creswell2022selection,imani2023mathprompter,taylor2022galactica}. Owing to training on vast and diverse textual data, LLMs have accumulated extensive prior knowledge, enabling them to fill the role traditionally occupied by domain experts in multiple contexts~\cite{Genomes2021,singhal2023large,wilcox2003role}, and facilitating automation across a spectrum of applications ranging from dialogue interpretation~\cite{finch2023leveraging} to automated program repair~\cite{fan2023automated}. More notably, by incorporating relevant task descriptions and limited exemplar data into their prompts, LLMs exhibit an aptitude for grasping contextual nuances and identifying salient features and instances within data~\cite{brown2020language,zhao2023survey}. This evidences their strong analytical skills and suggests their suitability for advanced feature engineering.

Recent approaches have sought to exploit the prior knowledge and inferential strengths of LLMs within tabular learning contexts. For example, background knowledge about demographic factors linked to disease risk can greatly augment disease prediction models. In such settings, contemporary studies articulate both task objectives and feature meanings in natural language, convert datasets into serial form, and subsequently input this structured information into an LLM~\cite{dinh2022lift,wang2023anypredict}. Downstream procedures commonly involve either in-context learning~\cite{nam2023semi} or fine-tuning techniques that alter only a subset of LLM parameters~\cite{dinh2022lift,hegselmann2023tabllm}. Leveraging the domain understanding inherently present in LLMs has proven to be effective, often surpassing conventional tabular learning approaches, especially in scenarios with very limited labeled examples~\cite{hegselmann2023tabllm}.

Prevailing methods that utilize LLMs for tabular data analysis encounter notable shortcomings. In end-to-end inference pipelines, each input requires at least one full LLM evaluation, resulting in significant computational overhead. Achieving high predictive accuracy often depends on fine-tuning the base LLM, which is impractical in scenarios where the model's training framework or fine-tuning API is inaccessible—a limitation that is increasingly relevant as state-of-the-art LLMs frequently restrict users to inference-only APIs~\cite{anil2023palm,openai2023gpt4}. Additionally, many of these approaches do not scale well when facing long textual prompts~\cite{liu2023lost}; as the dimensionality of tabular data increases, the required prompt length grows, rendering existing strategies either ineffective or prone to degraded performance. These obstacles present substantive barriers to deploying LLM-based tabular learners in practice.

Departing from end-to-end prediction paradigms, our approach repositions LLMs primarily as instruments for feature construction, thereby exploiting their inherent knowledge more strategically. We put forward a new framework, \textsf{FeatLLM}, that leverages in-context learning to decipher the “criteria” underpinning LLM decision-making. For example, within a medical diagnosis context, the LLM is capable of inferring and expressing rules or conditions—based on feature values—that indicate the presence of a disease. Such articulated rules can be employed to engineer new features that supplant existing variables in downstream learning tasks. This redirection confines LLM involvement to feature extraction, which reduces reliance on complex predictive models for subsequent inference and leads to lower inference time. Using LLMs’ in-context learning, these features can be derived without further training, simply by supplying suitable in-prompt exemplars. Moreover, drawing inspiration from ensemble techniques such as feature bagging~\cite{breiman2001random}, we can address the challenge posed by prohibitively long prompts due to numerous features.

\textsf{FeatLLM} decomposes the prompt design for automated feature construction into two main components: (1) comprehending the task and deducing the dependencies between the features and the target variables; and (2) leveraging insights from the earlier phase, along with limited labeled examples, to formulate discriminative rules that separate different classes. This stepwise logical procedure prompts the LLMs to disregard features that do not contribute meaningful information during rule generation. The extracted rules are operationalized on the dataset (treated as executable code), creating binary indicators for each rule that signal whether each data point satisfies the corresponding rule. These generated binary features are then employed by a simple classification model to estimate the probability of class membership. The stability and reliability of the results are further enhanced by an ensemble methodology, wherein feature generation is performed multiple times, combined with feature bagging. By tuning the number of features and the sample size used for each bag, \textsf{FeatLLM} alleviates issues caused by overly large prompts. Free from the need for parameter updates, this method operates without the burden of training costs and keeps inference efficient, broadening the practicality of LLM-based feature generation to diverse, real-world applications.

We assess \textsf{FeatLLM} across 13 varied tabular datasets under scenarios with limited labeled data, demonstrating consistently strong and stable results. We find that LLMs effectively combine pre-trained domain knowledge with dataset-specific evidence, resulting in the extraction of superior features. Our methodology surpasses leading few-shot learning approaches over several experimental contexts. Code and resources are publicly available at the anonymized GitHub repository~\texttt{https://github.com/Sungwon-Han/FeatLLM}.

\section{Related Work}

\subsection{Few-Shot Learning with Tabular Data} 

Few-shot learning methodologies have made it feasible to acquire broadly applicable representations from datasets while minimizing the need for labeled examples~\cite{chen2018closer}. Originally pioneered within the image domain, these approaches have since been extended to other data types such as text, audio, and, more recently, tabular data~\cite{majumder2022few,nam2023stunt,schick2021exploiting}. Relying on only a limited number of labeled instances introduces the potential for picking up misleading patterns~\cite{bartlett2021deep}. To address this, contemporary research has supplemented training with auxiliary data resources or injected domain-relevant priors to embed suitable inductive biases within models. Strategies include exploiting large collections of unlabeled data alongside well-crafted augmentation methods as in semi-supervised schemes~\cite{ucar2021subtab,yoon2020vime}, or utilizing unsupervised pre-training for model initialization independent of task specifics~\cite{somepalli2021saint}. Such unsupervised methods might leverage partial data masking or introduce noise and subsequently use reconstruction-based learning goals~\cite{arik2021tabnet,majmundar2022met}, or employ data augmentations to facilitate contrastive learning tasks~\cite{bahri2022scarf,chen2020simple}. Nonetheless, due to the diverse nature of tabular data, identifying effective augmentations that generalize well across different settings remains difficult, and these approaches often fall short for few-shot learning scenarios~\cite{nam2023stunt}.

More recent contributions advocate for model pre-training across a broad array of tabular datasets, even those unconnected to the final application, followed by leveraging transfer learning on the downstream task~\cite{zhang2023generative,levin2022transfer,zhu2023xtab}. For example, TransTab~\cite{wang2022transtab} trains transformer architectures to infer semantic connections between table columns, utilizing column names in addition to tabular values. The representations learned from multiple datasets and their attributes support generalization for zero-shot and few-shot tasks involving unfamiliar tables. In parallel, TabPFN~\cite{hollmann2022tabpfn} synthetically generates data of various distributions for the pre-training of a Bayesian neural network, which subsequently takes both the training and test data as input for inference on target tasks, all without further parameter updates. Such integration of prior information from wide-ranging datasets allows these methods to surpass traditional tree-based models, showing clear advantages in settings where labeled data is sparse.

\subsection{Language-Interfaced Tabular Learning} 

Incorporating large language models (LLMs)~\cite{anil2023palm,openai2023gpt4} has surfaced as an effective avenue for infusing external knowledge into tabular learning. Trained on expansive and varied textual datasets, LLMs exhibit substantial generalization capability, enabling them to address novel tasks across many fields~\cite{brown2020language}. Contemporary research has sought to exploit the breadth of knowledge encoded within LLMs for the benefit of tabular tasks. Typically, this involves converting tabular inputs into textual format and using these within the LLM's prompt interface~\cite{dinh2022lift,hegselmann2023tabllm,wang2023anypredict}. Prompts that blend data records, feature explanations, and descriptions of the intended task allow models to deduce meaningful relations between features and objectives, supporting effective inference. Approaches have diverged to either enhance LLMs for tabular scenarios via parameter-efficient adaptations such as LoRA~\cite{hu2021lora} or IA3~\cite{liu2022few}, or to utilize in-context learning by directly inserting a handful of example cases into prompts, eliminating the need for additional training~\cite{dinh2022lift,hegselmann2023tabllm,nam2023semi,slack2023tablet}.

Although the methods discussed above significantly advance few-shot tabular learning, their reliance on routing every prediction through an LLM introduces practical constraints, particularly in industrial environments. The present work seeks to move beyond this standard black-box paradigm that necessitates processing each sample end-to-end through the LLM. Instead, our focus shifts toward leveraging the LLM to derive explicit conditions or rules tied to each target class. \textsf{FeatLLM} translates these inferred rules into executable program code, which then produces novel features for downstream use. This shift allows subsequent predictions to be made independently of the LLM, capitalizing on in-context learning to distill useful rules from both pre-existing knowledge and the limited few-shot examples.

\section{Methods}
The core functionality of \textsf{FeatLLM} revolves around identifying and extracting the “rules,” that is, the defining conditions corresponding to each output class, from the limited set of available samples, as illustrated in Figure~\ref{fig:main_model}. After the LLM articulates these rules, they are programmatically parsed into code and subsequently utilized to construct binary features for a new sample, where each feature reflects the fulfillment (or lack thereof) of a particular rule. These newly generated binary indicators are then aggregated through a dot product with a set of non-negative, trainable weights to infer class probabilities. Training this architecture allows the model to assign significance to each feature and quantify its contribution to the class prediction (see Section~\ref{Sec:training}). This procedure is iteratively executed with bagging strategies, and the outcomes are aggregated through ensembling. The subsequent sections outline the overall problem setting and elaborate on the specifics of each component.

\vspace*{-0.12in}
\paragraph{Problem formulation.} Consider a tabular dataset comprising $N$ annotated instances, denoted as $\mathcal{D} = \{(\mathbf{x}^i, \mathbf{y}^i)\}_{i=1}^{N}$. Each sample $\mathbf{x}^i$ is a $d$-dimensional feature vector, and the collection includes $d$ input attributes accompanied by their respective natural language names $F = \{f_j\}_{j=1}^{d}$, and separately provided definitions. In a classification framework, each label $\mathbf{y}^i$ is assigned from a given set of categories. For the purpose of $k$-shot learning scenarios, only $k$ labeled examples, selected at random ($k<N$), are utilized during model training.

\subsection{Prompt Design for Extracting Rules}
\label{Sec:prompt_design}
To assist the LLM in deriving rules that reflect precise reasoning steps, we structure the interaction to emulate the analytical process of a domain expert tackling tabular problems. Our prompt is divided into three core sections (refer to Fig.~\ref{fig:prompt_for_rule} and Appendix~\ref{sec:example_rule}):

\vspace*{-0.12in}
\paragraph{Basic information description.} This segment outlines the primary details necessary for addressing the given task (see orange-highlighted section in Fig.~\ref{fig:prompt_for_rule}). It incorporates a clear statement of the learning objective with relevant label details, concise feature summaries, and example cases. The learning objective is articulated in a question format (such as, ``Does the myocardial infarction complications profile for this patient suggest chronic heart failure? Yes or no?". For more illustrations, see Appendix~\ref{sec:task_desc}). Feature summaries mention data types and can provide feature definitions, where applicable. For categorical attributes, sample category values may be included. Example demonstrations entail representing several labeled samples from the training set in text format, along with their correct labels. Serialization adheres to the conventions established in prior work~\cite{dinh2022lift,hegselmann2023tabllm}:

\begin{align}
&\text{Serialize}(\mathbf{x}^i,\mathbf{y}^i, F) = \nonumber \\ 
&``\ f_1\ \text{is}\ \mathbf{x}^i_1.\ ... \ f_d\  \text{is}\  \mathbf{x}^i_d.\ \ \text{Answer:}\  \mathbf{y}^i\ ”
\end{align}

\vspace*{-0.12in}
\paragraph{Reasoning instruction.} Our approach seeks to improve LLM reasoning by supplying targeted reasoning instructions (illustrated in blue in Fig.~\ref{fig:prompt_for_rule}). The reasoning instruction opens with a sentence akin to the chain-of-thought methodology~\cite{wei2022chain}. We then decompose the LLM's reasoning task into two distinct phases. During the initial phase, the model is prompted to deduce causal links or tendencies between features and the task description without referencing example demonstrations. This ensures reliance on general knowledge and common sense, enabling the model to maximize the use of its existing understanding relevant to the task. In the subsequent phase, the LLM considers provided examples in tandem with insights drawn from the first phase to articulate rules applicable to each class. This two-phase design steers the model away from forming misleading associations with non-informative columns and encourages prioritization of important features. To strike a balance between underfitting and overly verbose prompts, we fix the number of rules generated per class to 10.\footnote{Further discussion regarding rule quantity appears in Section~\ref{sec:analysis}.} Furthermore, the formulation of rules incorporates reference to the type of each feature from the basic information description section, reducing the risk of type mismatch within the rules.

\vspace*{-0.12in}
\paragraph{Response instruction.} To ensure that inferred rules are easily interpretable and practical, we instruct the LLM on the desired output structure by means of explicit response instructions (depicted in yellow in Fig.~\ref{fig:prompt_for_rule}). The prompt specifies detailed guidelines for how responses should be formatted for each answer class and delineates the structure that individual rules must adhere to. This guidance equips the LLM with criteria for choosing suitable output formats aligned with each feature type. In the absence of more restrictive instruction, the LLM retains flexibility to combine multiple rules using logical connectives such as AND or OR. A sample result is provided in Appendix~\ref{sec:outcome-rule}.

\subsection{Inferring Class Likelihood via Rules}
\label{Sec:training}

\paragraph{Parsing rules for feature generation.} The rules derived in the preceding phase serve as the basis for constructing new binary features. For every class, we generate features that take on values of either '0' or '1', depending on whether a sample complies with the set of class-specific rules. These features facilitate predictions regarding class membership likelihood for a given sample. Because the rules created by the LLM exist in natural language form, an appropriate parsing mechanism is necessary for automating the feature generation process. Although we attempt to standardize response and rule formatting to aid parsing, outputs from the LLM can nonetheless exhibit unpredictable syntactic fluctuations~\cite{kovavc2023large}. For example, rather than using symbols such as ``$>$” or ``$<$”, the LLM may generate textual alternatives like ``greater than” or ``less than,” complicating parsing procedures.

To mitigate issues arising from noisy textual representations, we avoid engineering elaborate code-based solutions. Instead, we enlist the LLM’s own capacities, since it demonstrates robust semantic interpretation skills and can produce simple code snippets in response to properly structured prompts (aided by its exposure to code in pretraining). Thus, our prompt to the LLM encapsulates the function definition, descriptions for both input and output, as well as the parsed rules. This prompt is then provided to the LLM, following examples given in Figures~\ref{fig:prompt_for_parsing} and ~\ref{sec:example_parsing}-\ref{sec:example_parse_outcome} in the Appendix. The Python \texttt{exec()} function is subsequently employed to execute the generated code for transforming the data.

\vspace*{-0.12in}
\paragraph{Inferring class likelihood.} Once these rule-based binary features for each class are generated, an intuitive approach for estimating the probability of class membership is to tally the number of satisfied rules per class (i.e., summing the class's binary features). Nevertheless, the predictive influence of each rule can differ, which necessitates learning their relative importance from labeled training data. To capture this, we apply a conventional machine learning (ML) model on the set of binary features for each class. For illustration, one can adopt a bias-free linear model. Denote by $\mathbf{z}_k^i \in \{0, 1\}^R$ the binary feature vector, representing the status of each of the $R$ class-$k$ rules for the $i$th sample. Here, $c$ denotes the number of classes. We then compute the class probabilities $\mathbf{p}^i$ as follows:

\begin{align}
\text{logit}_k^i &= \max(\mathbf{w}_k, 0) \cdot \mathbf{z}_k^i, \nonumber \\
\mathbf{p}^i &= \text{Softmax}({[\text{logit}_{1}^i, ..., \text{logit}_{c}^i]}). \label{eq:infer_prob}
\end{align}

In this approach, $\mathbf{w}_k$ denotes the adjustable parameters within the linear model, encoding the relevance of each corresponding feature. By constraining these weights to a positive domain, we guarantee that features generated by the LLM will not adversely affect the probability of class predictions during the training phase. Subsequently, the resulting logit vector is transformed into class probabilities via the softmax function. The learning of $\mathbf{w}_k$ is achieved by minimizing the cross-entropy loss, utilizing only a small set of training examples.

\vspace*{-0.12in}
\paragraph{Ensembling with bagging.} To form an ensemble for final prediction, we repeat the entire training and inference process multiple times, yielding several independently trained inference models. The aggregated prediction is generated by averaging the per-class probabilities $\mathbf{p}^i$ obtained for each class from all sets of trained weights $\{\mathbf{w}[t]\}_{t=1}^T$ across $T$ repeated trials, as described in Eq.~\ref{eq:infer_prob}.

To foster variability in rule generation, a sufficiently elevated temperature is employed during LLM inference. Moreover, the sequence of few-shot examples is shuffled since different orders can influence LLM outputs~\cite{lu2022fantastically}\footnote{Further examination of rule diversity appears in Appendix~\ref{sec:diversity}}. By embracing prediction diversity, we harness bagging to subsample features or data points in each trial. This is particularly advantageous when dealing with larger datasets or a more extensive feature set. The ensemble framework confers two principal benefits. First, the impact of erroneous rules produced by the LLM in certain trials is mitigated by the inclusion of correct rules from other iterations, thereby strengthening framework reliability. This aligns with the concept of self-consistency in LLMs~\cite{wang2022self}, as rules that repeatedly surface across different runs are likelier to be correct. Second, the ensemble mechanism addresses prompt length limitations intrinsic to LLMs. When tabular datasets surpass prompt capacity, applying bagging to limit input size ensures robust inference. While any single trial's rule set may insufficiently span the entire feature space, aggregating several weak learners developed from distinct trials compensates for such gaps, producing a model resilient to incomplete feature or instance coverage in single models.

\section{Experiments}
We assess the effectiveness of \textsf{FeatLLM} using a variety of tabular datasets, accompanied by extensive analyses to understand the behavior of our proposed framework. Owing to space limitations, further experimental results—including LLM variation, qualitative assessments, and additional studies—are deferred to the Appendix.

\subsection{Performance Evaluation}
\paragraph{Datasets.}
Our evaluation employs 13 different tabular datasets encompassing binary and multi-class classification tasks: (1) Adult~\cite{asuncion2007uci}, which predicts whether a person’s annual income surpasses \$50,000; (2) Bank~\cite{moro2014data}, for forecasting whether a client will subscribe to a term deposit; (3) Blood~\cite{yeh2009knowledge}, aiming to predict if a blood donor will make future donations; (4) Car~\cite{kadra2021well}, which assesses automobile quality; (5) Communities~\cite{misc_communities_and_crime_183}, which estimates the crime rate within a region; (6) Credit-g~\cite{kadra2021well}, for evaluating if a loan applicant is considered a favorable or unfavorable credit risk; (7) Diabetes\footnote{\texttt{kaggle.com/datasets/uciml/pima-indians-diabetes-database}}, which identifies individuals with diabetes; (8) Heart\footnote{\texttt{kaggle.com/datasets/fedesoriano/heart-failure-prediction}}, for detecting the likelihood of coronary artery disease; (9) Myocardial~\cite{misc_myocardial_infarction_complications_579}, designed to predict chronic heart failure status.

In addition, we incorporate recent and synthetically generated datasets that are not commonly evaluated and were absent from the pre-training corpus of relevant LLMs: (10) Cultivars, which predicts potential soybean cultivar grain yield\footnote{\texttt{archive.ics.uci.edu/dataset/913}}; (11) NHANES, for inferring an individual’s age category\footnote{\texttt{archive.ics.uci.edu/dataset/887}}; (12) Sequence-type, a dataset for identifying sequences as arithmetic, geometric, Fibonacci, or Collatz types; and (13) Solution-mix, which determines if the proportion concentration in a mix of four solutions exceeds 0.5. The first two datasets were newly added to the UCI Machine Learning Repository following September 2023, ensuring the LLM API (gpt-3.5-turbo-0613) leveraged in our work did not access these during pre-training. The remaining two are synthetic, precluding any memorization by the LLM, thus ensuring their validity for unbiased assessment.

A summary of dataset sizes and complexities can be found in Table~\ref{Tab:basic-desc}. All datasets provide distinct attribute names and clear descriptions. Notably, datasets such as `Communities' and `Myocardial' each possess more than 100 attributes, presenting significant difficulty due to LLM context window constraints.

\vspace*{-0.12in}
\paragraph{Baselines.}
We benchmark our approach against nine comparative methods. The initial trio consists of established tabular machine learning algorithms: (1) Logistic regression (LogReg), (2) XGBoost~\cite{chen2016xgboost}, and (3) RandomForest~\cite{ho1995random}. Subsequently, two methods that utilize unlabeled data include: (4) SCARF~\cite{bahri2022scarf} and (5) STUNT~\cite{nam2023stunt}; though, it should be noted that amassing large-scale unlabeled tabular datasets might not always be practical in operational contexts. As a recent innovation, (6) TabPFN~\cite{hollmann2022tabpfn} leverages pre-training on synthetically generated data from varied distributions. The remaining three baselines involve LLM-based methods: (7) In-context learning (In-context)~\cite{wei2022emergent}, (8) TABLET~\cite{slack2023tablet}, and (9) TabLLM~\cite{hegselmann2023tabllm}. The In-context learning baseline involves embedding a few examples in the LLM’s prompt without altering its parameters. TABLET adds additional reasoning cues, using external rule sets and prototypes via a classifier, to strengthen inference. TabLLM relies on parameter-efficient strategies (e.g., IA3~\cite{liu2022few}) to adapt the LLM with minimal additional data through few-shot learning.

\vspace*{-0.12in}
\paragraph{Implementation details.}
\textsf{FeatLLM} is implemented with GPT-3.5 serving as the foundational LLM, though its design permits flexibility in selecting alternative LLMs (refer to Appendix~\ref{sec:backbones} for experiments with other backbones). For LLM inference, we apply a temperature of 0.5 and maintain the top-p value at the default API setting of 1. The number of ensembles and the rule count for extraction are set at 20 and 10, respectively. Further discussion of the effects of hyper-parameter choices can be found in Figure~\ref{fig:hyperparameter2} and Appendix~\ref{sec:hyper-parameter}. The linear model utilizes the Adam optimizer with a learning rate of 0.01 and is trained across 200 epochs, employing $k$-fold cross-validation to identify the best-performing epoch. For reproducing baseline comparisons, GPT-3.5 is used in in-context learning baselines, while T0~\cite{sanh2022multitask} is chosen for TabLLM according to its original configuration. It is important to mention that TabLLM limits LLM use to models with open checkpoints, making GPT-3.5 unavailable for this setup. In the context of conventional machine learning models—specifically, LogReg, XGBoost, and RandomForest—parameter optimization is achieved via Grid-search combined with $k$-fold cross-validation. Here, $k$ is set to either 2 or 4, ensuring each class is represented in the training subset. Expanded implementation details are provided in Appendix~\ref{sec:baselines}.

\vspace*{-0.12in}
\paragraph{Results.}
Tables~\ref{tab:main1} and \ref{tab:main2} display comparative performance results across various datasets. Each experiment is conducted three times using different random splits of the training and test sets, with the reported AUC values reflecting both the mean and standard deviation. \textsf{FeatLLM} demonstrates strong results, frequently achieving either the highest or second-highest ranking among the baseline models. Remarkably, with only 4 shots, our method attains performance on par with standard tabular models that are trained with 32 shots, as highlighted in Fig.~\ref{fig:compares}. Although STUNT and TabPFN occasionally deliver substantial gains on individual datasets, their aggregate advantage over classic machine learning algorithms is modest. Furthermore, LLM-based models such as in-context learning, TABLET, and TabLLM encountered challenges processing tabular datasets with high dimensionality and failed to generate inferences in these settings. Our approach, integrating bagging and ensembling strategies, maintains robust performance even when handling datasets with many features. The results suggest that our ensemble and bagging methodology could be transferred to other LLM-driven techniques; however, such adaptations would necessitate doubling inference operations per sample, greatly amplifying computational demand and rendering the method impractical for routine use.

\subsection{Analysis \& Discussion}
\label{sec:analysis}
\paragraph{Ablation study.} To assess the relative impact of core framework components, we carry out ablation experiments examining the following aspects: (1) \textsf{-Tuning}, where the linear model's weight tuning is omitted and the logit is computed through a direct summation of binary features; (2) \textsf{-Ensemble}, where the ensembling mechanism is excluded; (3) \textsf{-Description}, where feature descriptions are removed; and (4) \textsf{-Reasoning}, where the Step-1 portion of the reasoning instruction is skipped during rule generation.

Table~\ref{tab:ablation} provides an overview of how the mean AUC changes for each ablation across all datasets. In every instance, eliminating a particular component leads to decreased performance. Notably, the positive effect of weight tuning intensifies as the number of shots rises, implying that with larger datasets, more precise rule importance estimations can be achieved and, thus, tuning yields greater benefit. Conversely, incorporating feature descriptions and reasoning instructions plays a more pronounced role with fewer shots, suggesting that capitalizing on LLMs' inherent prior knowledge is particularly advantageous in low-data regimes. Across all experimental settings, our ensemble approach consistently enhances model performance.

\vspace*{-0.12in}
\paragraph{Training \& inference time comparison.} We evaluate and contrast the training and inference durations across various approaches. All experiments are run on the Adult dataset using a training set with 16 examples. A single A100 GPU serves as the default computational resource, whereas TabLLM is distributed over four A100 GPUs to facilitate model parallelism. The comparative runtime data are reported in Table~\ref{tab:cost}. Importantly, \textsf{FeatLLM} achieves inference speeds close to those of traditional models, maintaining a notably low latency. Meanwhile, alternative LLM-based techniques—including In-context learning, TABLET, and TabLLM—incur significantly greater inference times. With respect to training time, \textsf{FeatLLM} remains competitive among LLM-driven methods, requiring only 30 API calls; this remains efficient in terms of budget and is amenable to parallel processing. 

\vspace*{-0.12in}
\paragraph{Quality of features generated by \textsf{FeatLLM}.}
To assess the practical usefulness of the rule-based features constructed by \textsf{FeatLLM}, we perform a targeted evaluation. We compute the average AUROC between each newly generated feature and the corresponding target label, then determine the proportion of features per dataset with AUROC values exceeding 0.5. Table~\ref{tab:quality} presents these findings. The results demonstrate that most of the constructed rules are indeed pertinent to the prediction task, offering valuable signal for solving the objective (see the ``Before tuning'' column). Furthermore, as a linear model is subsequently adjusted using data, rules exhibiting weak associations (those with learned weights less than a specified threshold) are pruned out automatically (see ``After tuning''). The threshold for feature retention is set to 0.1, based on inspection of the weight distributions.

\vspace*{-0.12in}
\paragraph{Effect of adding spurious correlations.} To evaluate the ability of different methods to filter out noisy, spurious correlations in few-shot learning settings, we performed an analysis using the Heart dataset as the primary task. Randomly selected columns from the Adult dataset, unrelated to the Heart task, were appended to the original features. Model performance was then assessed as the quantity of these extraneous columns increased incrementally. For an equitable assessment, each model was granted the same 8 labeled examples, omitting the use of unlabeled data and thus excluding methods such as SCARF and STUNT from the evaluation.

Figure~\ref{fig:spurious} illustrates how spurious correlations affect model performance. Among all evaluated approaches, \textsf{FeatLLM} demonstrated the smallest decline in accuracy as more irrelevant columns were introduced. This resilience appears to stem from our framework’s mechanism, which incorporates prior domain knowledge to guide the extraction of rules only from features relevant to the main task, thereby reducing reliance on irrelevant or noisy columns and achieving more stable outcomes. Empirically, we found that rules formulated from noisy, irrelevant features comprised just 1-2\% of the overall rule set. However, it is notable that \textsf{FeatLLM} is still susceptible to picking up spurious associations—particularly when appended columns originate from data distributions that are more closely aligned with the original task—potentially leading to errors in understanding causal structure (refer to Appendix~\ref{sec:spurious-appendix} for further discussion).

\vspace*{-0.12in}
\paragraph{Hyper-parameter analysis}
\textsf{FeatLLM} is governed by two primary hyper-parameters: the ensemble size and the count of rules per class retrieved in each LLM query. To understand how these settings influence overall effectiveness, we systematically evaluated their effects. Findings indicate that increasing the ensemble size enhances performance, though improvement plateaus around an ensemble size of 20 (see Fig.~\ref{fig:appendix-hyperparameter1} in Appendix). With respect to the number of rules per class, although no statistically significant variations were observed, we identify 10 rules as an optimal compromise between maintaining manageable prompt lengths and maximizing accuracy (refer to Fig.~\ref{fig:hyperparameter2}). We hypothesize this is because increasing the number of rules expands input dimensionality, thus supplying more information, but at the same time can induce overfitting, particularly in the low-shot scenario.

\section{Conclusion}
In this work, we introduce \textsf{FeatLLM}, which sets a new benchmark for large language model-based few-shot tabular learning. Unlike traditional paradigms where LLMs directly predict class labels per instance, \textsf{FeatLLM} leverages LLMs to synthesize decision criteria through automated feature engineering. By integrating rule extraction and bagged ensemble methods, \textsf{FeatLLM} attains low inference overhead, requires only access to LLM APIs (thus removing the need for additional training), and efficiently handles challenges posed by tabular data with large feature sets. Empirical results highlight its superior predictive capabilities, suggesting it is an effective and resource-efficient alternative for employing LLMs on practical tabular data problems.

\textsf{FeatLLM} specifically targets the low-shot learning context. Looking forward, we aim to expand the system’s feature engineering strategies for larger datasets, to incorporate a wider variety of feature construction methodologies beyond simple rules, and to enhance its interpretability. This direction holds promise for broadening applicability to more diverse use cases.  

\section{Broader Impact}
The framework presented, \textsf{FeatLLM}, seeks to capitalize on LLMs’ pretrained domain knowledge, thus enabling advancements over classical supervised tabular approaches, especially when data is scarce. By infusing data-driven learning with external knowledge, our contributions help address the common pitfall of overfitting in scenarios where labeled data is limited. In practical applications, such as finance and healthcare, where gathering labeled data is costly or otherwise constrained, \textsf{FeatLLM} offers meaningful benefits by tapping into the LLM’s relevant pretraining (often from vast bodies of finance- or healthcare-specific literature). Nonetheless, as reliance on LLM-derived priors grows, it becomes increasingly critical to study and mitigate societal biases or misinformation present in these models—particularly since such biases may unduly influence model outputs and, at times, may override evidence from the limited labeled data. This forms an important future avenue for making the technology both safer and more widely adopted.

\end{document}